\documentclass[
  degree=master,
  oneside,openany,
  AutoFakeBold=2.5
]{nuaathesis}
% 可用选项：
% degree=[bachelor|master|doctor]   必选，毕业论文学位
% type=[paper|design]   论文类型，本科必选，硕/博默认（强制）论文
% blankleft   可选，在启用 openright 时，去除空白页的页眉页脚，变成真正的白纸
% 其余选项将传给 ctexbook，比较有用的选项有：
% oneside|twoside       可选，单面或双面打印
% openany|openright     可选，在双面打印时，每章的第一页将打印在右侧
% AutoFakeBold          可选，貌似这个数值最接近 M$ Word 的粗体效果
% 推荐：双面打印时，使用 twoside,openright,blankleft；其他场合 oneside,openany

% 写作时，使用这个命令只渲染你想查看的部分，提升工作效率，定稿时注释掉整行
% \includeonly{chapter/00-cover,chapter/30-issues,}

\begin{document}

% 首先，导入论文的基础信息
%# -*- coding: utf-8-unix -*-

% 中文文档信息
\nuaaset{
  title = {\nuaathesis 示例文档}, % 论文题目
  author = {\nuaathesis~Group},   % 作者
  studentid = {012345678},  % 学号(本科)
  college = {\TeX 学院},    % 学院，或者(金城)院部
  major = {\LaTeX},     % 专业(本科)
  classid = {0123456},  % 班号(本科)
  advisors = {Donald Knuth\quad 大师},  % 指导教师
  libraryclassid = {TP371},       % 中图分类号(硕博)
  subjectclassid = {080605},      % 学科分类号(硕博)
  thesisid = {1028704 17-S036},   % 论文编号(硕博)
  majorsubject = {学位论文排版},  % 学科、专业(硕博)
  researchfield = {排版引用},     % 研究方向(硕博)
  % applydate = {二〇一七年七月}  % 默认当前日期
}

% 英文文档信息(本科可以不用写)
\nuaasetEn{
  college = {College of \TeX},
  title = {\nuaathesis~Tutorial},
  majorsubject = {Thesis Typesetting},
  author = {\nuaathesis~Group},
  advisors = {Prof.~Donald Knuth},
  % applydate = {July, 2017}
}

%% 中文摘要
\begin{abstract}
本文介绍如何使用\nuaathesis 文档类撰写南京航空航天大学学位论文。
\end{abstract}
\keywords{\TeX, \LaTeX, \nuaathesis}

%% 英文摘要
\begin{abstractEn}
A brief example of \nuaathesis.
\end{abstractEn}
\keywordsEn{\TeX, \LaTeX, \nuaathesis}


\makecover    % 封面
\makedeclare  % 承诺书

\frontmatter  % 启用页眉页脚
\makeabstract % 摘要
\nuaatableofcontents  % 正文目录
\nuaalistoffigurestables    % 图表清单

\mainmatter   % 以下正文

%# -*- coding: utf-8-unix -*-
\chapter{简介}\label{chap:intro}

这是南京航空航天大学(非官方)本科生学位论文\LaTeX 模板，当前版本是\version。本模板由\nuaathesis~Group共同开发，模板文档由Old Jack撰写。

本模板最早可以追溯到人人网上的一篇博客\footnote{\url{http://bit.ly/2rGOhf6}}，由黄大宁、邓欣珂、徐添豪、石坤四人共同开发完善，参考了当时东南大学的\seuthesix 模板；除此之外在Github上也可以找到一个repo\footnote{\url{https://github.com/nuaa803/nuaa-thesis}}，由Felix Ding、Jun Wang、Jackie Hou三位老师和Vevi Zhong同学共同维护，但是repo中的.cls和.sty文件是空文件。

回顾人人网的模板，没有直接提供nuaa.png和nuaa.bst文件，可以使用强制编译的方法生成文件，但是缺少左上角南航字样，参考文献格式也不符合标准。除此之外，旧模板使用了已经被放弃使用的CJK宏包，因此在编译\verb+\Unicode{}+命令时会出错，代码的阅读性和维护性也不如现在的ctex和xeCJK。由于上述原因，许多初次使用\LaTeX 和使用经验不多的同学，在一开始就放弃了使用旧版模板进行毕业设计的书写及排版。

基于南航无可用\LaTeX 学位模板可用的现状，\nuaathesis~Group基于旧\oldnuaathesis 模板、现东南大学的\seuthesix 模板和上海交通大学的SJTU Thesis模板，对模板进行了二次开发，基本实现了学士学位论文的模板。\nuaathesis~Group现有成员短期内不计划开发团队报告、硕士博士学位论文等其他文档的\LaTeX 模板，留给后续成员以及其他有需要、有能力的南航学子以后开发。

现在\nuaathesis 模板的代码托管在Github\footnote{\url{https://github.com/jackwzh/nuaathesis}}上，如有修改建议或者其他要求欢迎在Github上开issue或提pull request，\nuaathesis~Group会尽快回复，并酌情处理您的要求。

本模板基于Windows~10平台开发，使用MiKTeX v2.9发行版，所使用的宏包均跟进到最新版本。Linux平台由张一白使用\TeX Live测试，macOS平台由王成欣进行了测试，目前尚未出现任何问题。本模板尚未在Windows平台使用\CTeX / \TeX Live进行测试，如出现问题，请自行Google、Bing、Baidu搜索解决方法。学会使用搜索引擎、熟练阅读外文是一个学生最基本的能力，更是一个\LaTeX 使用者得以立足和前进的根本。

\nuaathesis~Group非常欢迎有其他南航的\LaTeX 使用者加入到本模板的开发与维护当中来，不断完善模板，为南航广大学子造福！

\section{模板使用}
\subsection{准备工作}
\begin{itemize}
  \item \TeX 发行版：Windows 系统推荐使用MiKTeX和\TeX Live这两种发行版，前者占用空间小，只在有宏包缺失情况下才进行下载，后者占用空间大，但基本无需担心宏包缺失。Linux系统(Arch系除外)推荐手动安装\TeX Live发行版，官方源中的TeXLive版本跟进较慢。macOS可以参考SJTU Thesis中的介绍。
  \item \TeX 知识：本说明文档提供\TeX 使用的例子，但不能解决所有的问题，因此使用前请自行学习\TeX~\&~\LaTeX 相关知识。
\end{itemize}

\subsection{模板编译}

\textbf{切记使用\XeLaTeX 进行编译。}使用\XeLaTeX 和\hologo{BibTeX}在文档中加入参考文献的流程可参考如下命令：

\iffalse
\begin{lstlisting}[basicstyle=\small\ttfamily, caption=手动逐次编译, numbers=none]
xelatex -no-pdf .tex文件名
biber --debug .tex文件名
xelatex .tex文件名
xelatex .tex文件名
\end{lstlisting}
\else
\begin{verbatim}
xelatex -no-pdf .tex文件名
biber --debug .tex文件名
xelatex .tex文件名
xelatex .tex文件名
\end{verbatim}
注：因为已知bug，临时禁用了 \texttt{lstlisting} 包。参见\url{https://github.com/CTeX-org/ctex-kit/issues/297}
\fi

使用\XeLaTeX 引擎编译可以直接通过各\LaTeX 编辑器实现，如：TeXworks，TeXmaker，TeXStudio，Emacs+插件，Atom+插件等等。biber命令需使用Windows的cmd/Power Shell、Linux和macOS下的bash实现。Windows平台可以自行编写简单的.bat批处理文件来实现。

使用biber需注意：\textbf{.bib文件内的文件记录必须在.tex文件中被引用，不引用的记录不要因为懒而不去除，否则将编译失败}

目录内容需要编译两次才能正常显示，原因推断为早期的电脑内存不够，所以将目录的生成分成了两步来进行。

\subsection{模板文件结构}
\begin{itemize}[noitemsep,topsep=0pt,parsep=0pt,partopsep=0pt]
  \item .tex文件：主文件，chapter下有各个章节的文件，强烈建议将文章模块化，方便调试与版本管理。
  \item .cls、.cfg文件：模板定义文件
  \item .bib文件：参考文献数据库文件
  \item figure文件夹：存放要插入的图片，其中nuaa.png不可删除
\end{itemize}

%# -*- coding: utf-8-unix -*-

\chapter{\LaTeX 排版用例}\label{chap:example}

本部分可参考SJTU Thesis模板\cite{SJTUThesis}的用例，源代码存放在tex文件夹下的examples.tex文件中，非常详尽。我们对其中几个部分有一些自己的见解，另外文献中也有一些没有提及的内容，写在本文档中。

\section{流程图}

对于不想花时间学习tizk宏包的同学，我们推荐在Power Point中绘制好流程图，然后导出pdf格式插入到文档中，学习成本较低，效果也非常好。

\section{表格}

SJTU Thesis中的表格介绍非常详尽，但其实有非常简单的从Excel生成\LaTeX 表格代码的方式，即Excel宏——excel2latex\footnote{\url{https://www.ctan.org/pkg/excel2latex?lang=en}}。此宏非常强大，可配合xcolor包生成有底色的表格。强烈推荐大家使用，提升撰写效率。

\begin{savenotes}
%% LaTeX通常不推荐在表格中插入脚注。但是，通过使用savenotes + makesavenoteenv 组合可在表格中插入脚注。请参考：https://en.wikibooks.org/wiki/LaTeX/Footnotes_and_Margin_Notes#Common_problems_and_workarounds

表\ref{tab:bicap}是张一白提供的一个\textbf{双语标题}和\textbf{在标题中使用脚注}的示例，对标题有特殊需求的同学可以参考此处的源代码：
\begin{table}[H]
  \centering
  \bicaption[chinese]{一个颇为标准的三线表格\footnote{这个例子来自\href{http://www.ctan.org/tex-archive/macros/latex/contrib/booktabs/booktabs.pdf}{《Publication quality tables in LATEX》}(booktabs宏包的文档)。这也是一个在表格中使用脚注的例子，请留意与threeparttable实现的效果有何不同。}}[english]{A standard three-line table\footnote{该表格演示了如何使用bicaption插入双语标题}}
  \label{tab:bicap}
  \begin{tabular}{@{}llr@{}} \toprule
    \multicolumn{2}{c}{Item} \\ \cmidrule(r){1-2}
    Animal & Description & Price (\$)\\ \midrule
    Gnat & per gram & 13.65 \\
    & each & 0.01 \\
    Gnu & stuffed & 92.50 \\
    Emu & stuffed & 33.33 \\
    Armadillo & frozen & 8.99 \\ \bottomrule
  \end{tabular}
\end{table}

\end{savenotes}

\LaTeX 中表格的使用体验比Word差很多，很遗憾这是不可避免的。除了使用上文提到的excel2latex宏之外，在线表格转换工具\footnote{http://www.tablesgenerator.com}也不失为一种高效的excel表格至\LaTeX 的转换方案。此外，也可以在word中打好表格，然后截图插入论文也不失为一种方案。

对于较长较大的表格，可以参考\LaTeX 笔记——lnotes2\footnote{\url{}http://dralpha.altervista.org/zh/tech/lnotes2.pdf}中的longtable（跨页表格）和sidewaystable（横向表格）等表格环境进行实现。另外可以使用\verb!p{2pt}!替代表格中的rcl，来控制表格每一列的宽度。

\section{图片表格指定位置插入}

图片和表格的插入默认是htbp四个选项，有时候这会让图片表格遍布整篇论文，可能会有同学非常反感这种情况，为了强制在当前位置插入图片，可以使用float宏包，然后使用H选项：\verb+\begin{figure}[H]+即可强制\TeX 在当前位置插入图片，从而避免正文和图片表格相距太远。

\section{多列图片}
由王成欣提供方案及示例代码。如需对两幅或多幅图片进行横向排版，建议使用subcaption包里的subfigure功能。效果如下：

\begin{figure}[H]
\centering
\begin{subfigure}{.45\textwidth}
  \centering
  \includegraphics[width=.9\linewidth]{nuaa-logo.pdf}
  \caption{左图}
  \label{fig:test_subfigure1}
\end{subfigure}
\begin{subfigure}{.45\textwidth}
  \centering
  \includegraphics[width=.9\linewidth]{nuaa-logo.pdf}
  \caption{右图}
  \label{fig:test_subfigure2}
\end{subfigure}
\caption{这是一个并列子图}
\label{fig:test_subfigure}
\end{figure}

请注意每行subfigure宽度的总和尽量不要超过一个\textbackslash textwidth，否则图像会自动折叠至下一行。

\chapter{一些来自 CQUThesis 的使用样例}

\section{字体命令}\label{txt:FreqCmd}
{\kaishu 玲珑骰子安红豆，入骨相思知不知。\hfill ——温庭筠}

{\fangsong 愿得一心人，白头不相离。\hfill ——卓文君}

% {\ifcsname youyuan\endcsname\youyuan\else[无 \texttt{youyuan} 字体。]\fi 去年今日此门中，人面桃花相映红。\hfill ——崔护}

{\heiti 入我相思门，知我相思苦。\hfill ——李白}

% {\ifcsname lishu\endcsname\lishu\else[无 \texttt{lishu} 字体。]\fi 此情可待成追忆？只是当时已惘然。\hfill ——李商隐}

{\songti 雨打梨花深闭门，忘了青春，误了青春。\hfill ——唐寅}

使用\texttt{textbf}和\texttt{textit}以及\texttt{underline}的效果分别如下：

这句话的\textbf{文字}分别\textit{使用}了三种命令来\underline{处理}。

The \textbf{words} in this sentences are \textit{processed} with three different \underline{cmd}.

（注：粗宋体可能会被替换成黑体，参见 \autoref{txt:issue:boldsun}）


\section{表格样本}

\subsection{基本表格}
\label{sec:basictable}

我们经常会在表格下方标注数据来源，或者对表格里面的条目进行解释。前面的脚注是一种不错的方法，如果不喜欢脚注，可以在表格后面写注释，比如\autoref{tab:tabexamp1}。
\begin{table}[htbp]
    \centering
    \bicaption{复杂表格示例}{A more structured table}
    \label{tab:tabexamp1}
    \begin{minipage}[t]{0.8\textwidth}
    \begin{tabularx}{\linewidth}{|l|X|X|X|X|}
        \hline
        \multirow{2}*{\diagbox[width=5em]{x}{y}} & \multicolumn{2}{c|}{First Half} & \multicolumn{2}{c|}{Second Half}\\\cline{2-5}
        & 1st Qtr &2nd Qtr&3rd Qtr&4th Qtr \\ \hline
        East$^{*}$ &   20.4&   27.4&   90&     20.4 \\
        West$^{**}$ &   30.6 &   38.6 &   34.6 &  31.6 \\ \hline
    \end{tabularx}\\[2pt]
    \footnotesize
    *：东部\\
    **：西部
    \end{minipage}
\end{table}

此外，\autoref{tab:tabexamp1} 同时还演示了另外两个功能：1）通过 \texttt{tabularx} 的\texttt{|X|} 扩展实现表格自动放大；2）通过命令 \texttt{diagbox} 在表头部分插入反斜线。

\autoref{tab:performance} 是一个很长的表格。

\begin{longtable}[c]{c*{6}{r}}
    \bicaption[实验数据]{实验数据，这个题注是双语的，而且十分的长，注意这在索引中的处理方式}[Data in experiment]{Data in experiment, and this is a really long long long long long long long long long long long long long long text.}\label{tab:performance}\\
    \toprule
    测试程序 & \multicolumn{1}{c}{正常运行} & \multicolumn{1}{c}{同步} & \multicolumn{1}{c}{检查点} & \multicolumn{1}{c}{卷回恢复}
    & \multicolumn{1}{c}{进程迁移} & \multicolumn{1}{c}{检查点} \\
    & \multicolumn{1}{c}{时间 (s)}& \multicolumn{1}{c}{时间 (s)}&
    \multicolumn{1}{c}{时间 (s)}& \multicolumn{1}{c}{时间 (s)}& \multicolumn{1}{c}{
        时间 (s)}&  文件 (KB) \\\midrule
    \endfirsthead
    \multicolumn{7}{c}{续表~\thetable\hskip1em 实验数据}\\
    \toprule
    测试程序 & \multicolumn{1}{c}{正常运行} & \multicolumn{1}{c}{同步} & \multicolumn{1}{c}{检查点} & \multicolumn{1}{c}{卷回恢复}
    & \multicolumn{1}{c}{进程迁移} & \multicolumn{1}{c}{检查点} \\
    & \multicolumn{1}{c}{时间 (s)}& \multicolumn{1}{c}{时间 (s)}&
    \multicolumn{1}{c}{时间 (s)}& \multicolumn{1}{c}{时间 (s)}& \multicolumn{1}{c}{
        时间 (s)}&  文件 (KB) \\\midrule
    \endhead
    \hline
    \multicolumn{7}{r}{续下页}
    \endfoot
    \endlastfoot
    CG.A.2 & 23.05 & 0.002 & 0.116 & 0.035 & 0.589 & 32491 \\
    CG.A.4 & 15.06 & 0.003 & 0.067 & 0.021 & 0.351 & 18211 \\
    CG.A.8 & 13.38 & 0.004 & 0.072 & 0.023 & 0.210 & 9890 \\
    CG.B.2 & 867.45 & 0.002 & 0.864 & 0.232 & 3.256 & 228562 \\
    CG.B.4 & 501.61 & 0.003 & 0.438 & 0.136 & 2.075 & 123862 \\
    CG.B.8 & 384.65 & 0.004 & 0.457 & 0.108 & 1.235 & 63777 \\
    MG.A.2 & 112.27 & 0.002 & 0.846 & 0.237 & 3.930 & 236473 \\
    MG.A.4 & 59.84 & 0.003 & 0.442 & 0.128 & 2.070 & 123875 \\
    MG.A.8 & 31.38 & 0.003 & 0.476 & 0.114 & 1.041 & 60627 \\
    MG.B.2 & 526.28 & 0.002 & 0.821 & 0.238 & 4.176 & 236635 \\
    MG.B.4 & 280.11 & 0.003 & 0.432 & 0.130 & 1.706 & 123793 \\
    MG.B.8 & 148.29 & 0.003 & 0.442 & 0.116 & 0.893 & 60600 \\
    LU.A.2 & 2116.54 & 0.002 & 0.110 & 0.030 & 0.532 & 28754 \\
    LU.A.4 & 1102.50 & 0.002 & 0.069 & 0.017 & 0.255 & 14915 \\
    LU.A.8 & 574.47 & 0.003 & 0.067 & 0.016 & 0.192 & 8655 \\
    LU.B.2 & 9712.87 & 0.002 & 0.357 & 0.104 & 1.734 & 101975 \\
    LU.B.4 & 4757.80 & 0.003 & 0.190 & 0.056 & 0.808 & 53522 \\
    LU.B.8 & 2444.05 & 0.004 & 0.222 & 0.057 & 0.548 & 30134 \\
    CG.B.2 & 867.45 & 0.002 & 0.864 & 0.232 & 3.256 & 228562 \\
    CG.B.4 & 501.61 & 0.003 & 0.438 & 0.136 & 2.075 & 123862 \\
    CG.B.8 & 384.65 & 0.004 & 0.457 & 0.108 & 1.235 & 63777 \\
    MG.A.2 & 112.27 & 0.002 & 0.846 & 0.237 & 3.930 & 236473 \\
    MG.A.4 & 59.84 & 0.003 & 0.442 & 0.128 & 2.070 & 123875 \\
    MG.A.8 & 31.38 & 0.003 & 0.476 & 0.114 & 1.041 & 60627 \\
    MG.B.2 & 526.28 & 0.002 & 0.821 & 0.238 & 4.176 & 236635 \\
    MG.B.4 & 280.11 & 0.003 & 0.432 & 0.130 & 1.706 & 123793 \\
    MG.B.8 & 148.29 & 0.003 & 0.442 & 0.116 & 0.893 & 60600 \\
    LU.A.2 & 2116.54 & 0.002 & 0.110 & 0.030 & 0.532 & 28754 \\
    LU.A.4 & 1102.50 & 0.002 & 0.069 & 0.017 & 0.255 & 14915 \\
    LU.A.8 & 574.47 & 0.003 & 0.067 & 0.016 & 0.192 & 8655 \\
    LU.B.2 & 9712.87 & 0.002 & 0.357 & 0.104 & 1.734 & 101975 \\
    LU.B.4 & 4757.80 & 0.003 & 0.190 & 0.056 & 0.808 & 53522 \\
    LU.B.8 & 2444.05 & 0.004 & 0.222 & 0.057 & 0.548 & 30134 \\
    EP.A.2 & 123.81 & 0.002 & 0.010 & 0.003 & 0.074 & 1834 \\
    EP.A.4 & 61.92 & 0.003 & 0.011 & 0.004 & 0.073 & 1743 \\
    EP.A.8 & 31.06 & 0.004 & 0.017 & 0.005 & 0.073 & 1661 \\
    EP.B.2 & 495.49 & 0.001 & 0.009 & 0.003 & 0.196 & 2011 \\
    EP.B.4 & 247.69 & 0.002 & 0.012 & 0.004 & 0.122 & 1663 \\
    EP.B.8 & 126.74 & 0.003 & 0.017 & 0.005 & 0.083 & 1656 \\
    \bottomrule
\end{longtable}

\section{定理环境}
\label{sec:theorem}

给大家演示一下各种和证明有关的环境：

\begin{assumption}
    假设以下数学方程成立：
    \begin{eqnarray}
    \label{eq:eqnxmp}
    c & = & a^2 - b^2\\
    & = & (a+b)(a-b)
    \end{eqnarray}
\end{assumption}

\begin{assumption}
    依然假设以下数学方程成立，注意整个系统是自动编号的：
    \begin{eqnarray}
    \label{eq:eqnxmp2}
    c & = & a^2 - b^2\\
    & = & (a+b)(a-b)
    \end{eqnarray}
\end{assumption}

\begin{cor}
    四川话配音的《猫和老鼠》是世界上最好看最好听最有趣的动画片。
    \begin{alignat}{3}
    V_i & =v_i - q_i v_j, & \qquad X_i & = x_i - q_i x_j,
    & \qquad U_i & = u_i,
    \qquad \text{for $i\ne j$;}\label{eq:B}\\
    V_j & = v_j, & \qquad X_j & = x_j,
    & \qquad U_j & u_j + \sum_{i\ne j} q_i u_i.
    \end{alignat}
\end{cor}

迢迢牵牛星，皎皎河汉女。
纤纤擢素手，札札弄机杼。
终日不成章，泣涕零如雨。
河汉清且浅，相去复几许。
盈盈一水间，脉脉不得语。

\begin{exmp}
    大家来看\autoref{ktc}。
    \begin{equation}
    \label{ktc}
    \left\{\begin{array}{l}
    \nabla f({\mbox{\boldmath $x$}}^*)-\sum\limits_{j=1}^p\lambda_j\nabla g_j({\mbox{\boldmath $x$}}^*)=0\\[0.3cm]
    \lambda_jg_j({\mbox{\boldmath $x$}}^*)=0,\quad j=1,2,\cdots,p\\[0.2cm]
    \lambda_j\ge 0,\quad j=1,2,\cdots,p.
    \end{array}\right.
    \end{equation}
\end{exmp}


\section{参考文献}
\label{sec:bib}
重庆大学的要求是参考文献以上标的形式标注于论述之后，就像这样：

研究表明\cite{r1}，早睡早起有益身体健康。

如果想同时引用多个文献\cite{r2,r3,r4,r6}，只需要在\verb|cite{}|中用逗号分开\textsf{citeKey}就好。
CQUThesis 同时提供正文模式的参考文献引用功能\texttt{inlinecite}，适用于以下情况：
文献\inlinecite{r6}表明，文献\inlinecite{r7,r8,r9}所述的情况是有理论依据的。

NUAAThesis 格式测试，学校的参考文献格式并不是 GB7714-2015，所以追加一些测试样例。
《要求》里列出的格式有：
\begin{enumerate}
  \item 连续出版物\cite{n11,n12}
  \item 专译集\cite{n21,n22}
  \item 论文集\cite{n31,n32}
  \item 学位论文\cite{n41,n42,n43}
  \item 专利\cite{n51,n52,n53}
  \item 技术标准\cite{n61,n62,n63}
\end{enumerate}

注：目前实现的格式仍然与《要求》有点差异，《要求》里论文集的编者、文集名、出版地是逗号分隔，
而目前是点号分隔；《要求》的学位论文用中文注明学位，目前没实现；在信息缺失的情况下，
《要求》貌似直接把对应字段省略，目前仍显示“XX不详”。

\chapter{已知问题}

这份模板还有很多欠缺的地方，有些问题虽然被发现了，
但目前还没找到满意的解决方法。本章记录模板还没解决的问题，希望用户能避开这些坑。
如果您有解决这些问题的思路，欢迎在 Github 上提交 issue 或 pull request。

\section{本模板的问题}

\subsection{英文 in 中文字体}

本模板部分英文的字体，与 Word 模板里不一样。

虽然要求里说，英文的字体要用 Times New Roman\textsuperscript{\textregistered}，
但 Word 模板里，不少地方的英文、数字的字体不是 Times New Roman\textsuperscript{\textregistered}，
而是旁边中文的字体。
比如正文的各级标题，还有本科封面，在 Word 模板里，这些西文的字体设置成“使用中文字体”。

\subsection{行间距}

行间距的问题由两方面造成：不知道 Word 是如何计算行间距；以及，不知道 \LaTeX 如何准确设置行间距。
这些不能准确计算的行间距目前只能通过测量来确定。

% 虽然理论上，《论文撰写要求》就应该完整地描述论文的格式。但实际上，……

网上能找到部分 Word 排版算法的解析\footnote{\url{https://www.zhihu.com/question/26397264/answer/48165229}}，
根据这些资料能计算出大部分《论文撰写要求》没有明确规定的格式，
但本科论文的页眉（页眉插入了一张图片）等地方，它们的排版算法与资料中的不一样。

\LaTeX 方面，貌似它会试图对齐所有页面的第一行？
在页首插入垂直空格时，会插入多余的空间。
还有，如果正文的各级标题出现在页首，那么它们的 \texttt{before} 段前间隔会被忽略掉。

\subsection{字号}

由于行间距的问题，请尽量避免使用规定字号以外的大小。

\LaTeX 貌似不能像 Word 那样设置成固定行间距，轻易修改全局字号、行间距可能会引发各种奇怪问题，
所以本模板只通过 ctex 的接口，设置了 \texttt{zihao} 和 \texttt{linespread}，
实现近似模板中的行间距效果。

\subsection{图表清单}

《要求》里并没有规定图表清单的格式，\LaTeX 里图清单和表清单是两个独立的列表，
所以每个列表里，章节之间会添加空行来提高可读性。
在两个列表合并到一起后，这些空行可能会反过来影响可读性，所以考虑将其去除……但是能力有限，没能去掉。
虽然参考了\footnote{\url{https://tex.stackexchange.com/a/331881}}，但貌似因为使用了某些宏包，
或者完全不是这种情况，所以没有成功。

\subsection{段落左侧缩进}

本科生正文要求段落左侧缩进2~字符，但\LaTeX 貌似有点麻烦，网上提到的方法大致可以分为两种：
一种把正文全都都加到一个环境里，另一种是每段前加一个命令来缩进。这两种都不是理想的解决方法，
参见讨论\footnote{\url{https://github.com/JackWzh/nuaathesis/issues/7}}。

\section{其他常见问题}

\subsection{宋体的粗体}\label{txt:issue:boldsun}

宋体的粗体可能会被换成黑体。
比如“\verb|右边是\textbf{粗宋体}|”，排版出来可能是“{右边是\heiti{粗宋体}}”。

这个问题源自很多宋体不带粗体，ctex 预置的6个字库中，只有 Fandol 和方正字库带有宋体的粗体。
但关键是 ctex 将其他字库中宋体的粗体，映射成黑体了。

解决方法有很多，一个方法是将代码写成“\verb|右边是\textbf{\songti 粗宋体}|”，
如果开启了 \texttt{AutoFakeBold}，就可以伪粗体的宋体了。

另一个方法是使用 Fandol 字库或方正字库，不过它们的粗宋体比 Word 论文模板里的要粗/重很多，
特别是作为标题使用时，差异会很大。

还有一些方法，比如修改字体映射之类，应该可行的样子，有兴趣的读者可以尝试研究一下。

\chapter{查重和其他注意事项}
\section{查重}

先说结论：{\large\textbf{知网完全支持pdf查重}}。

这个问题是鄙人整个毕设过程中最担心的问题之一，从知乎以及其他各种渠道搜索的结果并不一致；另外关于pdf查重具体检测哪些部分也是有很多种说法，现在根据鄙人论文的检测结果来说明一下几个需要注意的地方：

\begin{itemize}
  \item \textbf{页眉页脚：} pdf的眉页脚在论文查重检测范围内。如果担心会提升重复率，可以将页眉文字去掉（个人认为没必要）；
  \item \textbf{公式环境：} pdf中的公式在论文查重检测范围内。所以在编辑公式的时候，可以考虑不使用传统符号来编辑公式（物理公式符号不建议使用这种方法，各物理量的符号比较固定，老师可能会要求改正），以降低重复率，如参考文献中使用$\alpha$，可以改为$a$或$x$诸如此类；
  \item \textbf{表格环境：} 鄙人的论文中没有直接证据，但根据公式环境在查重检测范围内，鄙人推断表格的标题和内容很有可能也在范围内，所以建议大家不要直接摘抄实验数据和表格标题；
  \item \textbf{参考文献：} 鄙人在使用淘宝知网论文检测时，并未提交参考文献部分，学校不提供论文检测结果，所以目前没有直接证据证实参考文献是否在查重范围之内；
  \item \textbf{附录：} 鄙人的论文没有附录，情况不明。
\end{itemize}

鄙人的老师开始也要求上交word版论文，但是在鄙人的坚持下，最终上交了pdf版并成功通过查重。建议大家提前和指导老师打好招呼，最后提交pdf格式的论文。

\section{批注}
在论文撰写过程中，批注成了一个问题，鄙人的指导学姐并不是计算机专业出身，对\LaTeX 和基于Git的版本管理并不了解，所以沟通的途径就只有使用Adobe Acrobat等软件，对pdf文件本身进行批注，相比于word确实有些麻烦。

个人还是推荐使用Git\footnote{\url{https://git-scm.com/}}、Beyond Compare\footnote{\url{https://www.scootersoftware.com/}}等工具，辅以\LaTeX 本身的注释进行批注以及版本管理，非常清晰直观，操作也简单。

\section{毕业设计与毕业论文的区别}
这里特别对使用本模板的本科同学们做出提醒，请查看你们毕业设计基本信息中的毕设类别，共有两类：\textbf{毕业设计}和\textbf{毕业论文}。各位同学，你们\textbf{论文的封面和页眉中的内容应该与该类别相同}。因此在\verb!\documentclass[]{nuaathesis}!的选项中需要标明\textbf{Design}（毕业设计）或者\textbf{Paper}（毕业论文），使论文使用正确的封面和页眉。

除此之外该两类在最后论文装订时使用的并不是同一种封面纸，\textbf{毕业设计类的论文使用黄色的封面，毕业论文类的论文使用白色的封面}。在印刷厂/打印店打印时需提醒工作人员使用正确的封面纸张。

\section{单面打印\& 双面打印}
学校并没有规定论文打印的方式，考虑到部分同学有双面打印的需求，Gavin Lee 对twoside情况下的页脚进行了调整，奇数页页脚在右边，偶数页页脚在左边。可以在文档选项中使用oneside/twoside来切换单面打印和双面打印。

\section{封面打印\& 装订}
建议大家去印刷厂打印封面并装订。原因有下：
\begin{enumerate}
  \item 樱花广场打印店打印的封面并不标准，情况较复杂，总之是不标准的；
  \item 樱花广场打印店打印机并不稳定可靠，而且因为所有电脑都可以随意选择打印机，所以很容易出现打印错误，鄙人曾因员工操作失误以及机器故障被耽误2小时；
  \item 樱花广场打印店的档案袋储量较小，可能会用尽，而印刷厂不单独出售毕设档案袋，只能额外花钱买一整套封面来获取档案袋，存在浪费钱财的可能；
  \item 樱花广场打印店排队情况严重，因为有很多同学会在那里的电脑上修改他们的文档，从而影响了打印的效率。
\end{enumerate}

印刷厂虽远，但其质量是有保证的，封面也是标准的，另外因为距离远，排队现象相对较好，所以鄙人建议大家去印刷厂打印封面。

在印刷厂打印需要事先打好三个\textbf{A4纸}封面（论文封面、附件材料封面、工作材料归档封面），然后会使用你打印好的A4纸封面，复印到封面纸上，就得到了你的封面。

%# -*- coding: utf-8-unix -*-
\chapter{模板更新记录}
\label{chap:updatelog}

\textbf{2017年9月17日} v2.0发布，代码重构，模板结构重构，实现硕士博士模板，大量优化和修正，添加持续集成配置，添加大量代码示例。

\textbf{2017年6月22日} \nuaathesis 正式通过毕业设计审核，v1.0发布，增加毕业设计/毕业论文选项，并调整页眉；针对双面打印选项调整页脚；细节调整。

\textbf{2017年6月5日} v0.92发布，增加 biblatex 对 natbib 支持，如\verb!\citep!可以直接在行中引用编号， \verb!\citet!可以引用作者 (这里貌似仍然是个 bug, 理论上应该是引用题目，还没仔细研究。)； 添加subcaption和caption包，修复bicaption参数； 添加多列图片示例代码；多处细节调整。

\textbf{2017年5月12日} v0.91a发布，添加双语标题和标题中使用脚注用例；增加几个默认宏包，方便使用；部分细节修调整。

\textbf{2017年3月15日} v0.91发布，使用开源Fandol字体替代华文字体和思源雅黑字体。

\textbf{2017年3月14日} v0.9a发布，加入使脚注出现在页脚线下方的代码，加入模板更新记录。

\textbf{2017年3月14日} v0.9跨版发布，代码重构，模板基本实现，开始由Git进行版本控制，进入微调阶段。

\textbf{2013年6月4日} v0.3发布，加入对团队报告的支持，加入几个宏包，加一些预定义符号。

\textbf{2013年5月29日} v0.2发布，详情未知。

\textbf{2013年5月18日} v0.1发布，详情未知。

\textbf{2013年5月15日} 模板发布，版本号v0.0。


\backmatter   % 正文后无编号部分，
\bibliographystyle{gbt-7714-2015-numerical}   % 参考文献的样式
\bibliography{bib/sample}   % 参考文献

%# -*- coding: utf-8-unix -*-
\chapter[致谢]{致\hskip\ccwd{}谢}
首先，谨代所有以后使用本模板排版毕设的南航学子感谢本模板v2.0的代码贡献者——yzwduck，v2.0的工作基本均由他一人完成，让\nuaathesis 有了质的飞跃，向他表示最高的敬意！

其次要感谢使用、测试本模板v1.0版之前所有版本代码，并反馈问题、提pull request的几位co-authors，
是他们让\nuaathesis v1.0 得以如此美丽。
他们是：1613的张一白\footnote{\url{https://github.com/summershrimp}}和曾宪文\footnote{\url{https://github.com/RexSkz}}、0313的王成欣\footnote{\url{https://github.com/cvcore}}、0513的Gavin Lee\footnote{\url{https://github.com/gavinlee1994}}。
我从他们那学习到了很多新的\LaTeX 知识，也发现在南航有着和我一样喜欢\LaTeX 的朋友，而且他们都是比我更优秀的同学，我非常高兴能在大学的最后阶段认识他们。

最后特别感谢\oldnuaathesis 的四位作者：黄大宁、邓欣珂、徐添豪、石坤，\seuthesix 的开发维护团队，以及SJTU Thesis的开发维护团队，我们向他们的模板借鉴了很多的源码，没有这些前辈的工作，就没有今天的\nuaathesis，代表所有使用\nuaathesis 的南航学子感谢你们！

% \appendix     % 附录部分
\chapter{附\hskip\ccwd{}录}

\section{v0.9a后记——Old Jack 的吐槽}

\verb!\begin{轻松+愉快}!

Old Jack 他有点累......

Old Jack 两年前就开始关注南航毕设的\LaTeX 模板了，但是两年了还没有任何有实际意义的新动作，所以Old Jack 就亲自操刀制作了新的一版。虽然很多代码都是从其他模板中直接摘抄过来的，但是这也是\TeX 最普遍、最快捷的学习\&开发方法。一开始 Old Jack 也想造轮子，但是轮子真的不好造。

在制作过程中遇到了几个关键性的问题：
\begin{itemize}
  \item 前文提到的三种粗体
  \item nuaa.png源文件和页眉制作
  \item 英文字母、章节标题莫名其妙的加粗
  \item 脚注相对页脚线的位置
\end{itemize}

第一个问题 Old Jack 曾经用\TeX 中伪粗体(FakeBold)的方法实现过，但是效果并不好，而且当时受到最后一个问题的强烈影响，不得不使用其他字体来解决这个问题。

第二个问题 Old Jack 开始是使用官方模板中的图片，但是分辨率太低，效果很差。于是 Old Jack Google以图搜图找到了现在的这个文件的源文件，经过了一系列不可描述的操作后得到了现在的 nuaa.png 。页眉的制作也让 Old Jack 很头疼，论文要求论文到顶端和底端的距离分别为2.5cm和2.0cm，Old Jack 很naive的就给geometry设置了这个数值，但是效果和官方模板差了很多，于是 Old Jack 只好一点一点地调试，达到了近似官方模板的效果。页脚和官方模板有细微的区别，Old Jack 认为这无伤大雅，是要罗马数字和阿拉伯数字编号正确应该就可以了。

第三个问题是一个非常奇怪的问题。使用伪粗体时所有标题全都加粗了，非常难看，经过了代码重构和不停地调试解决了这个问题。在模板完成99\% 后发现最后致谢中的英文字体全都加粗了， Old Jack 几次审视代码和调试都没有解决。偶然间，Old Jack 将全部主要文件全部提取出来，放入另一个文件夹，然后重新编译就解决了这个问题！当然后来发现代码中确实有一个地方有小问题\textbf{可能}会影响，但是这不是上一次出错的原因。Old Jack 对于各位使用模板的南航学子以及其他可能会参考此模板的\TeX 爱好者提了一个建议：\textbf{任何语言，任何代码出现莫名其妙的问题时，换一个文件夹，改一下名字，重新跑一下，可能会得到意想不到的结果。}当然这不是万能的解决方法。

第四个问题就如第一章中脚注和页脚线的情况，感觉两条线很别扭。 Old Jack 犹豫了很久，最后没有采用将脚注放在页脚线下的方案，因为 Old Jack 觉得还是两条线的方案好看。对于想要将脚注放在页脚线下方的同学，可以在主文件中取消注释那段代码，来实现所需要的效果。

Old Jack 他完成了模板的再制作，但是他没有心气再写出一篇能够指导大家使用\LaTeX 的文档了(好吧，Old Jack 他承认懒是一部分因素)，望大家谅解 Old Jack。

\verb!\end{轻松+愉快}!

\section{v1.0后记}

Old Jack 非常高兴，因为他不是一个人在战斗。再次感谢张一白、王成欣、曾宪文、Gavin Lee等人的工作，没有他们，\nuaathesis 不会像现在这么美丽。

经过\nuaathesis~Group的努力和测试，\nuaathesis 迎来了v1.0版，也就是第一个正式发行版。一路走来也是有些坎坷，各种各样的小问题一直困扰着我们，其实v1.0 也还有着一些细小的问题尚未解决。不过Old Jack请大家放心，这些小问题不影响模板的使用。很多已经被我们解决的小问题比如页眉的大小位置，中英文字体是否正确，摘要的章节标题不能是加粗的宋体等等，老师可能不去管这些，甚至注意不到有什么区别。相比之下，重要的地方是：公式、图表的编号，图表和文本的位置，参考文献的格式等等才是老师关注的点。很多地方只是我们几个人为了追求和office模板尽可能接近，才不断地进行修改调整，也是有点讽刺。

写毕设论文的时候，Old Jack 不止一次看到隔壁室友调公式内容，Mathtype和Office装了卸，卸了装、调公式编号、调标题位置和大小、调首行缩进、调段间距等等等等，看着他们搞得焦头烂额的，Old Jack 都觉得心累。打印时也是这样，有太多的人在打印店不停地修改自己的论文，有因为office和wps不兼容修改的，有office版本不兼容修改的，有因为页眉页脚错误修改的等等。然而 Old Jack 他在写论文时从来没有担心过这些事情（当然，作为模板开发者 Old Jack 确实操心了很多，2333），他也第一次真正体会到了什么叫做专注于内容，真的挺轻松的（表格是例外）。

对于模板的推广，Old Jack觉得使用人数仍然不会太多，毕竟\LaTeX 的群众基础太小，除了8院，其他学院对公式的需求整体来讲并不迫切，Old Jack 猜测大部分知道、了解\LaTeX 的同学是通过数学建模竞赛这个途径才学习了\LaTeX ；同时因为涉及到学习新的程序语言，时间成本也较大，所以很多同学的学习意愿不高。不过\nuaathesis 的目标人群本来也不是全校所有学生，Old Jack 的思路，Old Jack 相信也是\nuaathesis~Group其他开发者的思路是：
\begin{enumerate}
  \item 为自己服务，这是\nuaathesis~Group开发模板的第一动力；
  \item 对已经掌握\LaTeX 基本语法的同学，\nuaathesis~Group为他们在毕业设计时能更轻松地撰写论文，提供平台和机会；
  \item 对准备学习\LaTeX 以及已经学习了一点\LaTeX 的同学，\nuaathesis~Group为他们提供学下去的动力和平台。
\end{enumerate}

即将毕业了，回首大学四年， Old Jack 做过疯狂的事情，也找到了一份看起来还可以的工作，只觉得还没对学校做过什么有用的事情，尽管 Old Jack 对学校其实并不是很有感情。完成了这个模板后，至少 Old Jack 可以减少一个遗憾，然后离开学校了。虽然这不是什么惊天动地的工作，但是至少 Old Jack 做了件他认为还算有意义的事情。Old Jack应该还会再维护\nuaathesis 一段时间，期待有后继者能够接过火炬，继续完善并推广\nuaathesis 。

想说的可能也就这么多了，Old Jack out！

\hfill 0813~王志浩，2017.6.24

\section{未知版本的后记}

也是两年前开始关注南航毕设的\LaTeX 模板了，但直到毕业前，都没能去静下心来学习\LaTeX。

现在差不多本科毕业一年，或者说，一年后要开写硕士学位论文了，决定这次一定要用\LaTeX。
本打算照着 CQUThesis 来造轮子的时候，逛纸飞机看到 \nuaathesis~v1.0 发布了。
非常激动、也很自愧，同样是经历了大学四年的人，我没能把这模板做出来。

于是马上把两年前为了模板而画的校名（矢量图）传了上去\footnote{\url{https://github.com/nuaatug/nuaathesis/commit/24fa82e4f136f3a713859335a395ecf364a342a2}}。

原本打算在 v1.0 版的基础上修改的，但因为行间距设置有问题，封面与 Word 模板也有点差异，
还要再加入硕/博士的模板，于是干脆改成 \texttt{Documented LaTeX Source (.dtx)}，
方便以后写模板的文档。

做模板过程中遇到的大问题，在于如何正确理解学校对论文格式的要求。
虽然有《本科毕业设计（论文）撰写格式要求》、《研究生学位论文撰写要求》，
但这些要求依然不够细致，因为那些要求都是假定你用 Word 来写论文的，要求里的内容是 Word 设置的操作方法，
所以还要先学习 Word 的排版算法。虽然这不是热门的资料，而且还有 CJK 独有的坑，
但幸好网上有高人已经把这问题基本解决了。

最后吐槽一下学校的 Word 模板，我觉得可以说，那个 Word 模板可能从最初做出来后，就基本没有变化。
那个“最初”或许可以追溯到上个世纪。很多编号的事情都要由手工来完成，比如说目录页码、
各级标题的编号、题注等。这些完全可以自动编号的工作，如果要手工做的话，【掀桌的emoji】。


\end{document}
